\section{A numerical second Iitaka construction}\label{num:second-base}

We record a numerical form of the second Iitaka construction for an
arbitrary nef contribution pulled back from an adjoint-positive base.
The original base may have nonnegative logarithmic Kodaira dimension;
its reduced inverse boundary is retained throughout the section
comparisons. The construction retains the dimension of the entire
second Iitaka base and, with empty base boundary, yields the ordinary
reduction in Section~\ref{ordinary:reductions}.

The argument separates the geometric construction from the numerical
step. First, the Kodaira-zero fibers determine a covering family on
the original base and hence a product diagram. Exact divisorial
orders then transfer sections to the second base. Finally, a uniform
volume estimate and interpolation show that this transfer retains
all of its dimensions. The nef divisor in the theorem is arbitrary;
its origin as a Hodge line is needed only in the subsequent ordinary
application.

\begin{theorem}[The dimension retained by interpolation]
\label{num:main}
Let $g:W\to Y$ and $h:W\to Z$ be surjective morphisms with connected
fibers between smooth projective complex varieties.  Let $B$ be an
effective rational SNC boundary, and let $D_Y$ be a reduced SNC
divisor, with the zero divisor allowed. Assume
\[
 B\geq(g^*D_Y)_{\red},\qquad
 \kappa(Y,K_Y+D_Y)\geq0.
\]
Let $A$ be a nef rational divisor on $Z$ such that $K_Z+aA$ is big
for all sufficiently large rational $a$.  Put
\[
 M=h^*A,\qquad L(c)=K_W+B+cM,\qquad d=\dim(W/Y),
\]
and suppose $L(1)$ is big over $Y$.
On smooth birational models, with the boundary rule of
Lemma~\ref{setup:log-modification}, there are morphisms
\[
 W\xrightarrow{q}U\xrightarrow{j}Z,
 \qquad U\xrightarrow{s}T_0,\qquad S\xrightarrow{p}T_0,
 \qquad \rho:S\to Y
\]
with the following properties:
\begin{enumerate}
\item $h=jq$, and $q$ is the relative Iitaka fibration of $K_W+B$
over $Z$.  Its geometric general fiber $G$ is integral and satisfies
$\kappa(G,K_G+B_G)=0$.
\item The map $g$ factors through $S$; $\rho$ is birational, and
$W\to(S\times_{T_0}U)_{\mathrm{main}}$ is dominant and generically
finite.  The generic fibers of $S\to T_0$ and $U\to T_0$ are
geometrically integral. With $D_S=(\rho^*D_Y)_{\red}$ made SNC,
one has $\kappa(S_t,K_{S_t}+D_S|_{S_t})=0$ for very general $t$.
\item One has
\begin{equation}\label{num:dimension-conclusion}
 \kappa(W,L(1))=\dim U=d+\dim T_0,
 \qquad \kappa(Y,K_Y+D_Y)\leq\dim T_0.
\end{equation}
\end{enumerate}
Rational linear equivalence in the definition of $M$ is sufficient.
All section comparisons are with the fixed initial source model.
\end{theorem}

A divisor is big over $Y$ when its restriction to the geometric
generic fiber of $g$ is big. All boundaries in the theorem have
coefficients in $[0,1]$. A point has Kodaira dimension zero, so the
statement includes zero-dimensional bases and fibers.

Setting $D_Y=0$ recovers the ordinary numerical construction, including
$D_S=0$ and $\kappa(S_t)=0$. In the logarithmic setting the inverse
boundary condition is needed both to contract the base pluriform and
to apply the model-pair comparisons.

\subsection{A covering family with Kodaira dimension zero}

The first task is to determine the dimension of the family of images
of a Kodaira-zero fiber. A logarithmic pluriform on $Y$ gives both
nonnegativity for each moving image and the differential-rank
calculation for the family.

Fix $0\ne\alpha\in H^0(Y,\ell(K_Y+D_Y))$, replacing it by a tensor
power whenever needed so that $\ell$ clears the boundary denominators
in the following comparisons. We will use two consequences of
pulling it back to a family covering $Y$.

First, a general member $S_b$ of any covering family of integral
subvarieties of $Y$, resolved with reduced inverse boundary $D_{S_b}$,
has $\kappa(S_b,K_{S_b}+D_{S_b})\geq0$.
Indeed choose $\dim Y-\dim S_b$ parameter directions spanning its
generic normal space.  The family over a transverse parameter slice
has dimension $\dim Y$ and maps generically finitely to $Y$.
Pull back $\alpha$ to a resolution of this family and its inverse
boundary, and use adjunction on the general fiber. This produces a
nonzero logarithmic pluriform on a resolution of $S_b$. The assertion
also holds after a finite cover with its reduced inverse boundary.
These image boundaries consist only of the reduced inverse image
of $D_Y$; exceptional divisors outside that inverse image are not
automatically added. Logarithmic pullback is regular away from that
inverse image, and this convention ensures compatibility with the
allowed boundary on the source. The construction uses the family of
\emph{images}; its domains need not map generically finitely to $Y$.

\begin{lemma}[Rank of the image family]\label{num:cycle-rank}
Let $q:V\to U$ and $g:V\to Y$ be surjective morphisms between
smooth connected projective complex varieties.  Suppose the general
fiber $G$ of $q$ is geometrically integral, $B$ is an effective
rational boundary, $(V,B)$ is log smooth over a
dense open of $U$, $D_Y$ is reduced SNC,
$B\geq(g^*D_Y)_{\red}$, $\kappa(Y,K_Y+D_Y)\geq0$,
$\kappa(G,K_G+B_G)=0$, and
$G\to g(G)$ is generically finite.  Then the family of reduced
image cycles $g(G)$ has dimension $\dim Y-\dim G$.
\end{lemma}

\begin{proof}
Put $r=\dim G$ and $k=\dim Y-r$.  On a smooth parameter open,
the image family is flat.  At a general $y\in S=g(G)$ the normal
evaluation map
\[
 e_y:T_uU\longrightarrow N_{S/Y,y}
\]
has rank $k$, since the family covers $Y$.  Its nonzero maximal
minors are obtained from transverse $k$-dimensional parameter slices.
Pullback of $\alpha$ to each sliced family and adjunction produce
sections of $\ell(K_G+B_G)$.  Their local coefficients are the
$\ell$-th powers of the minors, multiplied by the same tangential
Jacobian and the coefficient of $\alpha$. The reduced inverse-boundary
inclusion makes every logarithmic pole allowed by $B_G$, proving
global regularity in the indicated logarithmic bundle.
Changing lifts of parameter vectors
by vectors tangent to $G$ does not change their wedge determinant.

Every nonzero section space of a divisor with Iitaka dimension zero
is one-dimensional.  Ratios of the resulting sections are therefore
constant on $G$.  The common factors cancel, so all nonzero minor
ratios are constant as well: their $\ell$-th powers are constant,
and a rational function on an integral complex variety with a
constant positive power is itself constant.  Thus the Pl\"ucker point of the
rank-$k$ quotient $e_y$, and hence $\ker(e_y)$, is independent of
general $y$.  A parameter vector in that kernel induces an embedded
deformation of $S$ zero at its generic point.  Such a deformation is
a homomorphism from the conormal sheaf to $\mathcal O_S$; since $S$
is integral, it is detected at the generic point and must be zero.
The common kernel is therefore the kernel of the image-cycle map.
In characteristic zero its generic differential rank is its image
dimension, namely $k$.
\end{proof}

\subsection{Construction of the product}

We first verify that a relative Iitaka fibration exists over $Z$.
After constructing it, the preceding rank calculation will identify
the parameter space of its images in $Y$.

For a general fiber $W_z$ of $h$, let $S_z$ be its resolved image
in $Y$, with reduced inverse boundary $D_{S_z}$. The covering-family
observation gives $\kappa(S_z,K_{S_z}+D_{S_z})\geq0$.
The fibers of $W_z\to S_z$ have big log canonical divisor.  To see
this restriction carefully, write a multiple of $L(1)$, over a
dense open of $Y$, as a relatively ample divisor plus an effective
divisor.  General intersections with $h$-fibers move in $g$-fibers
and avoid the fixed effective support.  The ample part remains ample
on them; $M$ is trivial there; and adjunction after resolving the
combined map identifies the restriction of $K_W+B$ with their log
canonical divisor. The reduced inverse boundary on $S_z$ pulls back
into $B_{W_z}$. The lc general-type bound following
Equation~\eqref{num:KP} now gives
\[
 \kappa(W_z,K_{W_z}+B_{W_z})\geq0.
\]
Take the relative Iitaka fibration $W\xrightarrow qU\xrightarrow jZ$.
The relative Iitaka theorem \cite[Section~6, Remark~1]{IitakaOriginal} gives
\begin{equation}\label{num:iitaka-dimension}
 \kappa(G,K_G+B_G)=0,
 \qquad\dim U=\dim Z+\kappa(W_z,K_{W_z}+B_{W_z}).
\end{equation}
One applies it to the full log divisor on every chosen model:
Lemma~\ref{setup:log-modification} identifies the relative section
sheaves as well as the global section spaces.

If $G\to g(G)$ had positive-dimensional fibers, the same ample-plus-
effective restriction would make their log canonical divisors big.
The image has nonnegative logarithmic Kodaira dimension by the covering-family
observation. The source boundary on $G$ contains the reduced inverse
boundary of that image, so the lc general-type bound would then contradict
$\kappa(G,K_G+B_G)=0$.  Hence $G\to g(G)$ is generically finite.

Let $T_0$ be a smooth model of the relative algebraic closure in
$\C(U)$ of the image-cycle field.  Pull the universal reduced image
family to $T_0$ and resolve its dominating component to obtain
$S\to T_0$.  The regular extension $\C(U)/\C(T_0)$ ensures that
$U\to T_0$ has geometrically integral generic fiber.  The generic
image member is geometrically integral too: after this extension it
is dominated by the geometrically integral generic fiber of $q$.
Thus the fiber product has one component dominating $T_0$.  The maps
fit into
\begin{equation}\label{num:product}
\begin{tikzcd}[column sep=large]
 W\arrow[r]\arrow[d,"q"']&
 (S\times_{T_0}U)_{\mathrm{main}}\arrow[r]\arrow[d]&
 S\arrow[r,"\rho"]\arrow[d,"p"']&Y\\
 U\arrow[r,equal]&U\arrow[r,"s"']&T_0&
\end{tikzcd}
\end{equation}
By Lemma~\ref{num:cycle-rank}, $\dim T_0=\dim Y-\dim G$.
Both the top-left map and $\rho$ are consequently generically finite,
and $\dim U-\dim T_0=d$.  Moreover
$\C(Y)\subseteq\C(S)\subseteq\C(W)$, with the first extension
finite.  Connectedness of $g$ makes $\C(Y)$ algebraically closed in
$\C(W)$, so $\rho$ is birational.

Put $D_S=(\rho^*D_Y)_{\red}$ and resolve this support. The source
boundary contains $(g_S^*D_S)_{\red}$, since both reduced inverse
boundaries are supported over $D_Y$. The pullback of $\alpha$ to
$S$ restricts, by adjunction, to a nonzero logarithmic pluriform
on $S_t$. Logarithmic pullback to $G$, which generically finitely
dominates $S_t$ with the required inverse boundary, injects its
systems into those of $K_G+B_G$. Therefore
$\kappa(S_t,K_{S_t}+D_S|_{S_t})=0$.
Because $\rho$ is birational and $D_S$ is exactly the reduced
inverse boundary, logarithmic pullback and pushdown at every original
prime identify the divisible section spaces of $K_Y+D_Y$ and
$K_S+D_S$. Equivalently, the latter spaces lie between the pulled-back
spaces and the identical spaces obtained by adding every reduced
exceptional divisor, as in Lemma~\ref{prelim:sections}.
The elementary upper
addition bound gives
\begin{equation}\label{num:base-bound}
 \kappa(Y,K_Y+D_Y)
 =\kappa(S,K_S+D_S)\leq\dim T_0.
\end{equation}
It remains to prove that $L(1)$ retains all $\dim U$ directions.

\subsection{The divisor measuring descending forms}

The product diagram determines a rank-one family of forms along
$q$. We now record exactly how much a coefficient from $\C(U)$ may
have a pole while the associated form stays regular on the original
source. These coefficients, rather than a limiting assertion about
Kodaira dimensions, will give the transfer used in the last step.

Fix the original model $W_0$ whose sections are to be recovered.
Apply Lemma~\ref{ro:extraction} to $q$, and preserve its conclusion
relative to $W_0$ under the further modifications below.  Thus every
divisor of a source model with codimension-two image on $U$ is
exceptional over $W_0$.

Write $M_U=j^*A$. Take the $\ell$-th tensor power of a rational
relative top form on $U/T_0$, wedge it tensorwise with the logarithmic
$\ell$-pluriform $\alpha$ on $S$, and pull back through the generically
finite map in Equation~\eqref{num:product}. In each common degree $m$
divisible by $\ell$, the $(m/\ell)$-th power of this combined form
spans the entire fiber section space, since $\kappa(G,K_G+B_G)=0$.
After trivializing the bundles pulled back from $U$, division of a
degree-$m$ section by this power gives a ratio constant on the geometric
generic $q$-fiber; connectedness of $q$ puts that ratio in $\C(U)$.

For a prime $P\subset U$, let $\tau_P$ generate $\ell K_{U/T_0}$
at its generic point and let $\eta_P$ be its combined rational
$\ell$-pluriform on $W$.  Define
\begin{equation}\label{num:order-divisor}
 c_P=\min_{E\mapsto P}
 \frac{\ell^{-1}\operatorname{ord}_E(\eta_P)+b_E}
 {\operatorname{ord}_E(q^*P)},\qquad
 C=\sum_Pc_PP,\qquad P_U(c)=K_{U/T_0}+C+cM_U.
\end{equation}
Here $b_E$ is the coefficient of $B$ and the order of $\eta_P$ is
measured in $\ell K_W$.  Only finitely many coefficients are nonzero:
outside the divisors of one fixed rational frame, its combined form,
$B$, and the divisorial nonsmooth locus, every order is zero.

Fix a rational value of $c$ and take common degrees divisible by
$\ell$ and the indices of all divisors involved. Division by the
fixed rank-one form, and multiplication by that form in the reverse
direction, identify the full divisible section systems:
\begin{equation}\label{num:descent}
 H^0(W,mL(c))\xrightarrow{\ \sim\ }H^0(U,mP_U(c)).
\end{equation}
For the forward map, a rational coefficient of order $a$ at $P$ has order
\[
 a\operatorname{ord}_E(q^*P)+(m/\ell)\operatorname{ord}_E(\eta_P)+mb_E
\]
upstairs, so its allowed orders are exactly $a\geq-mc_P$.
Conversely a section of $mP_U(c)$ satisfies every divisorial test on
$W_0$: divisors dominating $U$ are handled by the generic rank-one
form, divisors over base primes by this minimum, and all other
divisors are exceptional over $W_0$.  Normality then gives a section
on $W_0$.  The converse lands on the fixed original model $W_0$, whose systems
agree with those on the current model by
Lemma~\ref{setup:log-modification}. Consequently these are mutually
inverse comparisons of the full divisible section systems. Tensor
products use powers of the same generating form, so multiplication
and section ratios are preserved, as are the dimensions of their
rational-map images.

\begin{lemma}[The horizontal and vertical orders]\label{num:orders}
On suitable fixed models, write $C=C_h+C_v$ relative to $U\to T_0$.
The divisor $C_h$ is an SNC boundary over the generic point of $T_0$,
$C_v\geq-R_s$, and
\begin{equation}\label{num:order-lower}
 P_U(c)\geq K_{\mathcal V_s}+C_h+cM_U.
\end{equation}
For very general $t$, the same minimum construction for
$q_t:W_t\to U_t$, using the adjunction form on $S_t$, gives
$C_h|_{U_t}$.  In particular, for every rational divisor $Q$ on $U_t$,
\begin{equation}\label{num:fiber-descent}
 H^0(W_t,m(K_{W_t}+B_{W_t}+q_t^*Q))
 \lhook\joinrel\longrightarrow
 H^0(U_t,m(K_{U_t}+C_h|_{U_t}+Q)).
\end{equation}
\end{lemma}

\begin{proof}
First check specialization.  Fix rational frames and include their
divisors, the combined form, $B$, and all divisorial nonsmooth data in
one finite support.  Record separately their codimension-two images
on $U$.  At very general $t$, vertical members are absent and
horizontal strata meet the fiber transversely.  Their coefficients,
orders, and multiplicities are unchanged.  Every minimizing divisor
$E\to P$ remains surjective after base change, so every component of
$P_t$ receives an unchanged minimizing value.  The complements of the
order-computation loci have smaller fiber dimension and contain no
such component.  Thus the minimum remains $c_P$.

A listed divisor dominating $U$ cannot create an extra vertical test:
over the geometric generic point of $T_0$, its components are
conjugate, and one dominates the geometrically integral $U_t$, hence
all do.  This spreads after shrinking the base.  A codimension-two
image stays of codimension at least two, or is avoided.  Outside the
finite support the frames are units and the generic divisor fibers
are smooth and reduced.  The fiber minimum is therefore exactly
$C_h|_{U_t}$.  Division by the rank-one fiber form proves
Equation~\eqref{num:fiber-descent}; no extension of a fiber section
to $W$ is asserted.

For a horizontal $P$, work on such a fiber.  A regular top form on
$U_t$ wedged with the logarithmic pluriform on $S_t$ pulls back
with only the poles allowed by $B_{W_t}$,
so $c_P\geq0$.  In the finite normalization of $S_t\times U_t$,
choose a divisor over $S_t\times P_t$.  Its ramification index $e$
gives the value $(e-1+b_E)/e\leq1$, because the form from $S_t$ is
a unit at its generic point.  Thus $c_P\leq1$.

For a vertical prime $P$, work at its discrete valuation ring, with
uniformizer $x$, and choose a local volume form $\theta$ on $T_0$.
If $r_P$ is the coefficient of $R_s$, write
$s^*\theta=x^{r_P}\beta$ with $\beta$ primitive.  The saturated
subspace defined by this decomposable form is a direct summand over
the discrete valuation ring.  There is consequently a regular
complementary $d$-form $\lambda$ with $\beta\wedge\lambda$ a unit
top form.  As a relative form, $\lambda$ represents $x^{r_P}\tau$,
where $\tau$ generates $K_{U/T_0}$.  Wedging $q^*\lambda$ with the
form from $S$ shows that $(q^*x)^{\ell r_P}\eta_P$ is allowed by
$B$: its only possible logarithmic poles lie in the reduced inverse
boundary already contained in $B$. Therefore $c_P\geq-r_P$. This argument also applies when the
image of $P$ has codimension greater than one on $T_0$.

Finally resolve the finite support on $U$, resolve the source map,
and recompute the minimum.  At an old prime, old minimizing
valuations persist and pullback of a log-regular form stays regular
with the boundary rule of Lemma~\ref{setup:log-modification}.
Thus old coefficients do not change.  New horizontal coefficients
satisfy the same product calculation, giving SNC horizontal support
and coefficients in $[0,1]$.  Equation~\eqref{num:saturation} gives
Equation~\eqref{num:order-lower}.
\end{proof}

\subsection{Complete systems and finite group actions}

Return to Theorem~\ref{num:main}. Fix a rational $c\in(0,1)$ close
enough to one that $L(c)$ is still big over $Y$, and fix the models
in Lemma~\ref{num:orders}. The next objective is to prove that, for
very general $t\in T_0$,
\begin{equation}\label{num:fiber-big}
 D_t:=K_{U_t}+C_h|_{U_t}+cM_U|_{U_t}\quad\text{is big}.
\end{equation}
When $d=0$, this is the zero-dimensional assertion. For $d>0$ we
will bound the volumes of small ample perturbations of $D_t$ below
by one positive constant. This bound must be independent of the
perturbation: full Iitaka dimension for each positive perturbation
alone would not prove bigness of $D_t$ in the limit.

The zero-excess case of Proposition~\ref{num:model-pair} will make
the auxiliary general-type model pairs constant after finite
extensions. The extensions may vary with the perturbation. The next
two results retain complete model systems in all divisible degrees
and bound the loss when taking invariants by the effective action
on the section image. The first supplies one fixed normalization
and one fixed open set for every required degree of each such model.

\begin{proposition}[Complete systems in an isotrivial family]
\label{num:full-systems}
Let $p:H\to I$ be a surjective morphism of smooth integral
projective varieties over an algebraically closed field of
characteristic zero. Let $B$ be an effective rational SNC boundary on
$H$, let $D_I$ be a reduced SNC divisor on $I$, and assume
$B\geq(p^*D_I)_{\red}$ and $\kappa(I,K_I+D_I)\geq0$.
Put $r=\dim H-\dim I$. Choose a geometric generic fiber
component after adjoining the Stein field. Suppose it is klt of log
general type and its log canonical model pair is isotrivial; denote
its fixed model by $(F_c,\Delta_c)$. There are a finite Galois base
extension, the normalization $\tau:H^a\to H$ in the corresponding
compositum component, and an integer $m_0>0$ such that
\begin{equation}\label{num:full-system-injection}
 H^0(F_c,m(K_{F_c}+\Delta_c))
 \lhook\joinrel\longrightarrow H^0(H^a,m\tau^*(K_H+B))
 \qquad(m\in m_0\N).
\end{equation}
On one fixed dense open of a smooth altered base, restriction to a
general fiber gives the pulled-back complete system of this model,
times a common nonzero factor. Both the open and $m_0$ work for all
these degrees. The isotriviality hypothesis holds in particular if
$\kappa(H,K_H+B)\leq r$.
\end{proposition}

\begin{proof}
Write $C=B_{\mathrm{hor}}$, and let
$V=R_p+B_{\mathrm{vert}}-p^*D_I\geq0$ as in
Equation~\eqref{num:decomposition}. Prescribe the finite extension
identifying the chosen generic model
with $(F_c,\Delta_c)$, including its Stein field, in
Lemma~\ref{num:alteration}. A model section defines a relative rational
pluriform on $H'$. At horizontal primes it is allowed by the generic
canonical-ring comparison. At a vertical prime tested by
Lemma~\ref{num:alteration}, work over $R=\mathcal O_{I',D}$ and compare
with $(Z,\Delta_Z)=(F_c,\Delta_c)\times\operatorname{Spec}R$.
Its reduced special fiber $Z_0$ makes $(Z,\Delta_Z+Z_0)$ log canonical.
On a common graph resolution, the discrepancy $a_E$ of the tested
valuation satisfies
\[
 a_E-\operatorname{ord}_E(t)\geq-1,
 \qquad \operatorname{ord}_E(t)=1,
 \qquad a_E\geq0.
\]
Use the same base canonical divisor on both sides of the comparison,
so its pullbacks cancel and these are also relative discrepancies.
In a degree clearing the model index, the zero divisor of a model
section is effective Cartier. Its pullback is effective, and the
remaining relative-form order is $ma_E\geq0$. Thus every model section
is regular at every tested prime, including those centered on the
boundary or singular locus of the fixed model.

Lemma~\ref{num:alteration} gives sections of
$m\tau^*(K_{\mathcal V_p}+C)$. Fix
$0\ne\alpha\in H^0(I,\ell(K_I+D_I))$. Multiplication by its
pulled-back $m/\ell$-th power and by the section of $m\tau^*V$ gives
Equation~\eqref{num:full-system-injection}, using
Equation~\eqref{num:decomposition}. All factors are independent of
the chosen model section and nonzero on a general fiber.

To fix one open for every degree, use finite generation of the generic
log canonical ring \cite{BCHM}. Choose a divisible index $m_1$ for
which its Veronese is generated in degree one and
$A_c=m_1(K_{F_c}+\Delta_c)$ is basepoint-free Cartier. Write
$L_\eta=K_{H'_\eta}+C'_\eta$ on the generic fiber of $p'$.
Resolve the base ideal of $|m_1L_\eta|$ and its model map on one
common graph, with maps $u$ to $H'_\eta$ and $v$ to the fixed model.
The free and fixed parts give
\[
 u^*(m_1L_\eta)=v^*A_c+E_{\mathrm{fix}},
 \qquad E_{\mathrm{fix}}\geq0.
\]
Generation in degree one implies, for every $a\geq1$,
\[
 |a u^*(m_1L_\eta)|=v^*|aA_c|+aE_{\mathrm{fix}}.
\]
Spread these maps and this divisor identity over one dense open of
$I'$. Remove the images of its finitely many vertical terms and the
zeros of the base multiplier, as well as $p'(\Supp\rho^*V)$.
Increase $m_0$ to a common multiple of
$m_1$ and the earlier indices. Every restricted comparison is then a
tensor power of the same free-and-fixed-part identity, so this open
and index work for all $m\in m_0\N$. Injectivity is already visible
on the chosen generic fiber. The last assertion follows from
Proposition~\ref{num:model-pair}.
\end{proof}

The alteration making a general-type model constant can have arbitrarily
large degree.  The next observation avoids dividing section growth by
that degree.  It bounds only the action on the image of the system.

\begin{lemma}[A degree bound for invariant growth]\label{num:invariants}
Let $D$ be a rational Cartier divisor on a normal projective complex
variety $H$, with $\kappa(H,D)=d>0$.  Let $H\dashrightarrow V$ be dominant,
where $\dim V=d$, and let $\pi:H'\to H$ be a finite normal Galois
cover.  Suppose that in every sufficiently divisible degree $m$:
\begin{enumerate}
\item the image field of $|m\pi^*D|$ contains $\C(V)$;
\item one fixed integral subvariety $F'\subset H'$ dominates that
image and maps generically finitely to $V$ by the induced map, with
degree at most $e_0$;
\item a subspace of the restriction image of
$H^0(H',m\pi^*D)$ on $F'$ is, up to a common rational factor
nonzero at the generic point of $F'$, the pullback under a dominant
rational map from $F'$ of the complete degree-$m$ system of a fixed
ample rational divisor on a $d$-dimensional projective model, whose
volume is at least $v>0$.
\end{enumerate}
Then
\begin{equation}\label{num:invariant-growth}
 \limsup_n\frac{d!\,h^0(H,nD)}{n^d}\geq\frac{v}{e_0}.
\end{equation}
The cover, the ample model, and the degree index may vary in
applications; $v$ and $e_0$ alone control this bound.
\end{lemma}

\begin{proof}
Choose one sufficiently divisible $m$ for which the ample-model
system is an embedding, and let $X_m$ be the complete-system image.
Finite pullback preserves Iitaka dimension, so $\dim X_m=d$.
Because $F'$ dominates $X_m$, restriction of the global sections
defining $X_m$ is injective on their image: a section whose
restriction vanishes gives a hyperplane containing the dense image
of $F'$, and hence vanishes on $X_m$. The subspace in (3) therefore
lifts to a linear subsystem on $X_m$. Its rational linear projection
has as its image precisely the embedded ample model. A general
linear section of the projected image lifts to a linear section of
$X_m$; the points in a general finite fiber are counted with their
positive multiplicities. Hence
$\deg X_m$ is at least the product of the projection degree and
the degree of that model. In particular,
\[
 \deg X_m\geq vm^d.
\]
The pullback linearization gives an action of the Galois group
$\Gamma$ on the complete system and on $X_m$. It fixes the
embedded field $\C(V)$, since this field comes from $H$.  The
tower
\[
 \C(V)\subseteq\C(X_m)\subseteq\C(F'),
 \qquad[\C(F'): \C(V)]\leq e_0
\]
shows that its effective image on $X_m$ has order $e\leq e_0$.

Let $R$ be the homogeneous coordinate ring of $X_m$.  Evaluation
embeds $R_k$ equivariantly into $H^0(H',km\pi^*D)$.  Its invariants
descend to $H^0(H,kmD)$, since
$(\pi_*\mathcal O_{H'})^\Gamma=\mathcal O_H$.
Choose a positive integer $r$ divisible by the orders of the
characters with which every stabilizer in the \emph{full} group
$\Gamma$ acts on the fibers of $\mathcal O_{X_m}(1)$; for example,
$|\Gamma|$ suffices. Thus the line bundle $\mathcal O_{X_m}(r)$
with its given linearization descends to an ample line bundle $A_m$
on the finite quotient $X_m/\Gamma$. The quotient has generic
degree $e$, the order of the effective action, even when the kernel
acts nontrivially by scalars on the original line bundle. It follows
that $A_m^d=r^d\deg(X_m)/e$.

For sufficiently large $n$, the coordinate-ring piece $R_n$ agrees
with $H^0(X_m,\mathcal O_{X_m}(n))$: this is the restriction
sequence in projective space and Serre vanishing for the ideal
sheaf of $X_m$. Finite quotient descent identifies the invariant
sections in degree $rk$ with $H^0(X_m/\Gamma,A_m^k)$.
The Hilbert polynomial therefore gives
\[
 \dim R_{rk}^{\Gamma}
 =\frac{r^d\deg(X_m)}{e\,d!}k^d+O(k^{d-1}).
\]
Normalize in degree $mrk$ and let $k$ tend to infinity. Cancelling
$r^d$ gives the lower bound
\[
 \frac{\deg(X_m)}{em^d}\geq\frac{v}{e_0}.
\]
Thus even an arbitrarily large scalar kernel affects only the
divisibility index, not the claimed lower bound.
\end{proof}

\subsection{Uniform volume on a fiber of the second base}

We now prove Equation~\eqref{num:fiber-big} with $c$ and the models
fixed as above. Choose a very general $t\in T_0$. The case $d=0$
has already been noted, so assume $d>0$ and abbreviate
\[
 H=W_t,\qquad I=S_t,\qquad V=U_t,
 \qquad e_0=\deg(H\to I\times V).
\]
Put $D_I=D_S|_{S_t}$. We have
\[
 \kappa(I,K_I+D_I)=0,\qquad
 B_H\geq(p_H^*D_I)_{\red},\qquad \dim V=d,
\]
where $p_H:H\to I$ is the induced morphism, and $L(c)|_H$ is
big over $I$.
The fiber order test in Equation~\eqref{num:fiber-descent} bounds
the Iitaka dimension of every divisor $L(c)|_H+q_t^*Q$ by $d$.

Fix an ample rational divisor $H_V$ on $V$ and, for rational
$\delta>0$, put
\[
 N_\delta=L(c)|_H+\delta q_t^*H_V.
\]
Before choosing $\delta$, decrease the coefficients of $B_H$
horizontal over $I$ slightly to obtain $B_H^-$, leaving all vertical
coefficients fixed. Thus the coefficient-one inverse image of $D_I$
is retained. Choose the decrease
so small that $K_H+B_H^-+cM|_H$ remains big on every geometric
generic fiber component over $I$.  Its finitely many component
volumes have a fixed positive lower bound $v$.  Adding the nef
divisor $\delta q_t^*H_V$ cannot lower this bound.

The divisor $cM_U|_V+\delta H_V$ is ample.  A sufficiently divisible
general member of its pulled-back system gives an snc rational
boundary $\Theta_\delta$ with small coefficients.  Thus
\[
 E_\delta:=K_H+B_H^-+\Theta_\delta
 \sim_\Q N_\delta-(B_H-B_H^-)
\]
has klt general fibers over $I$, with volume at least $v$.
The pair still contains the reduced inverse image of $D_I$, and its
logarithmic base has Kodaira dimension zero. Equation~\eqref{num:KP},
the inclusion of these systems in those of
$N_\delta$, and the rank-one upper bound imply
\[
 d\leq\kappa(H,E_\delta)
 \leq\kappa(H,N_\delta)\leq d.
\]
Their general-type model pairs consequently have variation zero.
Proposition~\ref{num:full-systems} supplies a finite normal Galois
cover $\tau:H^a\to H$ induced by an extension of $\C(I)$, with
complete ample-model subseries of volume at least $v$.  Multiplying
by the section of $B_H-B_H^-$ places them in the systems of
$\tau^*N_\delta$.  This common factor is nonzero on a geometric
general fiber.  The extension and the model may depend on $\delta$.

We check the two remaining conditions of Lemma~\ref{num:invariants}.
The same argument for $\delta/2$ gives a nonzero section of a
multiple of $N_{\delta/2}$.  Its powers multiplied by the systems of
$(\delta/2)q_t^*H_V$ have ratios generating $\C(V)$ in sufficiently
divisible degrees.  Hence the complete-system image field contains
$\C(V)$, both downstairs and after finite pullback.

Next choose a smooth comparison model $\widehat H\to H^a$ mapping
to the altered smooth base $I'$ and to its fixed canonical model.
Use the fixed open and the fixed index in
Proposition~\ref{num:full-systems}; these arise from one graph and
one divisor identity, rather than degree-dependent open sets.
Choose a very general fiber component $\widehat F$ and let
$F'\subset H^a$ be its reduced image.  The birational map
$\widehat H\to H^a$ is birational on this component, since a
geometric general fiber component is not contained in its exceptional
locus.  Its restricted model system has image dimension $d$ and therefore
dominates the $d$-dimensional complete-system image in every required
degree.  This fixes a single $F'$ for the application of the lemma.

Finally, $[\C(F'): \C(V)]\leq e_0$, independently of the altered
base. To see the precise bound, write $K_H=\C(H)$,
$K_I=\C(I)$, and $K_{I'}=\C(I')$. For the selected compositum
component of the base change, its degree over $I'\times V$ is
\[
 e'=[K_HK_{I'}:K_{I'}(V)]\leq
       [K_H:K_I(V)]=e_0.
\]
Shrink one dense open of $I'$ so that the corresponding finite
map has this generic degree on its fibers. In characteristic zero
the geometric generic fiber is reduced, and the degrees of its
components over $V$ sum to $e'$. The same holds for sufficiently
general fibers after this shrinking. The chosen $F'$ is birational
to one such component, so its degree is at most $e'$. Intersect this
open with the open already fixed for the model systems and choose
$F'$ there. The common multiplying sections are nonzero at its
generic point, and the same $F'$ works in every required degree.
This argument includes any preliminary Stein extension and never
divides by $[K_{I'}:K_I]$.

Lemma~\ref{num:invariants} now gives
\[
 \limsup_n\frac{d!\,h^0(H,nN_\delta)}{n^d}\geq v/e_0.
\]
The fiber order test injects these spaces into those of
$n(D_t+\delta H_V)$.  Since $\dim V=d$, we obtain the ordinary
volume bound
\[
 \operatorname{vol}_V(D_t+\delta H_V)\geq v/e_0
 \qquad(\delta\in\Q_{>0}).
\]
Continuity of volume \cite[Section~2.2.C]{Lazarsfeld} gives
$\operatorname{vol}_V(D_t)\geq v/e_0>0$, proving
Equation~\eqref{num:fiber-big}.

\subsection{Interpolation on the second base}

Only one positivity input remains: relative bigness gives a
pseudoeffective saturated relative divisor, whereas a big divisor
from the base supplies all base directions. We state both forms of
weak positivity before interpolating between them.

\begin{lemma}[Two consequences of weak positivity]\label{num:weak}
Let $p:H\to I$ be a surjective morphism between smooth connected
projective complex varieties.
\begin{enumerate}[label=\textup{(\roman*)}]
\item If $\Delta$ is a rational SNC boundary, $Q$ is a big rational divisor
on $I$, and the generic fiber is connected, then
\[
 \kappa(H,K_{H/I}+\Delta+p^*Q)
 \geq\dim I+\kappa(H_i,K_{H_i}+\Delta_i).
\]
\item If $C$ is a rational horizontal boundary log smooth over the generic
point of $I$, $P$ is a nef real divisor, and $K_H+C+P$ is big over $I$, then
$K_{\mathcal V_p}+C+P$ is pseudoeffective.
\end{enumerate}
\end{lemma}

\begin{proof}
Apply Lemma~\ref{num:alteration}, including the Stein field when
needed. For a log smooth pair with connected fibers, twisted weak
positivity says that
\[
 \mathcal E_m=p'_*\mathcal O_{H'}(m(K_{H'/I'}+C'))
\]
is weakly positive in divisible degrees
\cite[Theorem~1.1]{FujinoWeakPositivity}. Its reflexive symmetric
powers are generically generated after arbitrarily small positive
ample twists.

Here is the extension step needed when $p'$ is not equidimensional.
Since $\mathcal E_m$ is torsion free, it is locally free on an open
$J\subset I'$ containing every codimension-one point. On $J$, the
reflexive symmetric power agrees with the ordinary symmetric power,
and multiplication of sections defines the regular evaluation map
\[
 \operatorname{Sym}^b\mathcal E_m\longrightarrow
 p'_*\mathcal O_{H'}(bm(K_{H'/I'}+C')).
\]
Thus a global section of the reflexive symmetric power with a base
line-bundle twist evaluates to a rational pluriform regular over
$J$. Every possible divisorial pole lies over $I'\setminus J$.
Lemma~\ref{num:alteration} makes such a divisor exceptional over the
original $H$; its order test at all other primes and normality give
a section on $H^a$. This argument is the replacement for a global
reflexivity assertion about the direct image.

For (i), the assertion is automatic if the fiber Iitaka dimension is
$-\infty$.  Otherwise discard the vertical boundary and write the pulled-back big
divisor as $Q'\sim_\Q2\eta A+E$, with $A$ ample, $\eta>0$ rational,
and $E\geq0$.  Choose a fiber degree $m$ realizing the fiber Iitaka
dimension $k$.  Weak positivity with the first $\eta A$ supplies
systems whose restrictions include all products of those degree-$m$
fiber sections, so their fiber image has dimension $k$.  Multiplying
by the sections of the second $\eta A$ and by the fixed section of
$E$ makes the image field contain the base field.  The resulting
image has dimension $\dim I+k$.  Lemma~\ref{num:alteration} descends
these systems to the finite normalization of $H$; adding $R_p$ and
the discarded vertical boundary embeds them in the required systems.
Finite pullback preserves Iitaka dimension. This proves the
big-base comparison by the weak-positivity method of
\cite[Proposition~9.1]{FujinoWeakPositivity}; the divisorial tests
above supply the extension step that equidimensionality supplies
in that formulation.

For (ii), first assume $P$ is rational, and take all perturbation
parameters positive and rational. Fix ample divisors $A'$ on $H'$
and $A_I$ on $I'$.
Represent $\rho^*P+\delta A'$ by a general small-coefficient SNC
boundary.  Its addition preserves the big generic log divisor.
The same weak-positivity argument, with a twist $\epsilon A_I$,
produces a nonzero section of a sufficiently divisible positive
multiple of
\[
 K_{H'/I'}+C'+\rho^*P+\delta A'+\epsilon p'^*A_I
\]
regular at every prime tested over $H$. Push its divisor to $H$: any
negative components are exceptional over $H$, so the pushforward is
effective. Divide by $e=\deg\rho$. The tested lattice comparison
and the norm for principal divisors show that
\[
 K_{\mathcal V_p}+C+P+
 \frac{\delta}{e}\rho_*A'+\frac{\epsilon}{e}\rho_*p'^*A_I
\]
is pseudoeffective.  These two error classes are fixed.  Let
$\delta,\epsilon$ decrease to zero and use closedness of the
pseudoeffective cone.

For a nef real divisor $P$, choose ample rational divisor classes
$Q_j$ tending to $[P]$ in $N^1(H)_{\R}$, by approximating the ample
classes $P+\epsilon_j A_H$ with $\epsilon_j\downarrow0$ and $A_H$
ample. Their restrictions tend to $[P|_{H_i}]$ on each geometric
generic fiber component. Bigness is open, so $K_H+C+Q_j$ is big over
$I$ for all sufficiently large $j$. The rational case makes
$K_{\mathcal V_p}+C+Q_j$ pseudoeffective. Closedness of the
pseudoeffective cone gives the assertion for $P$.
\end{proof}
We now have every ingredient on the same fixed model $U$.
By Lemma~\ref{num:weak}(ii) and Equation~\eqref{num:fiber-big},
$K_{\mathcal V_s}+C_h+cM_U$ is pseudoeffective.  Thus
Lemma~\ref{num:orders} makes $P_U(c)$ pseudoeffective.
Choose rational $a>1$ with $K_Z+aA$ big.  Applying
Lemma~\ref{num:weak}(i) to $h$ gives
\[
 \kappa(W,L(a))
 \geq\dim Z+\kappa(W_z,K_{W_z}+B_{W_z})=\dim U.
\]
The section identification in Equation~\eqref{num:descent} makes
$P_U(a)$ big.
Both divisors lie on the same affine line, and
\[
 P_U(1)=\frac{a-1}{a-c}P_U(c)+\frac{1-c}{a-c}P_U(a).
\]
The two coefficients are positive, so $P_U(1)$ is big. The inverse
comparison in Equation~\eqref{num:descent} transfers its
systems to $L(1)$ on the original $W_0$, preserving their rational
map dimensions.  Therefore $\kappa(W,L(1))\geq\dim U$.
The opposite inequality is elementary upper addition for $q$, since
$M|_G=0$ and $\kappa(G,K_G+B_G)=0$.  Together with
Equation~\eqref{num:base-bound} and the product dimension identity,
this completes the proof of Theorem~\ref{num:main}.
